% =====================================================================
% Two paper subsections for the IMTS forecasting experiment:
%   (a) how we used S5 and RoMAE (baseline-setup paragraph),
%   (b) interpretation of why each baseline succeeds or fails (analysis).
% Designed to slot into the experiments section of the main paper.
%
% Citation keys assumed:
%   smith2023simplified   — Smith, Warrington, Linderman (S5, ICLR 2023)
%   zivanovic2025rotary   — Zivanovic et al. (RoMAE, arXiv:2505.20535)
%   zhang2024irregular    — Zhang et al. (T-PatchGNN, ICML 2024)
%   gu2023mamba           — Gu, Dao (Mamba, 2023)
% Cross-references assumed:
%   tab:imts_real_results — main results table (results_table_phase5.tex)
%   sec:method            — our method section
% =====================================================================

\subsection{Baseline Setups: S5 and RoMAE}
\label{sec:baselines_setup}

We benchmark our model against two architecturally distinct foundation
baselines, in addition to the seventeen specialized IMTS baselines reported by
\citet{zhang2024irregular}.

\paragraph{S5 (Smith et al., ICLR 2023).}
We use the paper-native pendulum configuration ($d_{\text{model}}{=}128$,
state dim $P{=}256$, $n_{\text{layers}}{=}6$, AdamW with weight decay $0.05$
excluded from the SSM-structured parameters $\Lambda$, $\log\Delta$, $B$, $C$,
$D$), giving 1.40M parameters. Per-step continuous-time $\Delta t$ is threaded
through S5's \texttt{step\_scale} argument, identical to the paper's pendulum
recipe (Section~5.4 of \citet{smith2023simplified}). For multivariate IMTS,
each variate is processed independently by the same parameter-shared S5
stack---channel-independent at the variate level, MIMO at the
$d_{\text{model}}$ level. This is a forced design choice: variates in raw IMTS
do not share observation timestamps, so true MIMO across the variate axis
(stacking $V$ values per timestep) would require pre-alignment to a common
grid, which S5 itself does not provide. Per-variate channel-independence is
the only honest pure-S5 application to raw IMTS.

\paragraph{RoMAE (Zivanovic et al., 2025).}
We use the configuration $d_{\text{enc}}{=}288$, $\text{depth}_{\text{enc}}{=}7$,
$d_{\text{dec}}{=}180$, $\text{depth}_{\text{dec}}{=}2$, $p_{\text{RoPE}}{=}0.75$,
giving 7.80M parameters (interpolated between the paper's tiny and small
variants to match a TSFM-scale budget). Each observation becomes one flat
token with continuous position vector $[\text{timestamp},\ \text{variate\_id}]$,
exactly as recommended in Section~4.2 of \citet{zivanovic2025rotary} for
multivariate IMTS, with the learnable $[\text{CLS}]$ token included to break
the relative-position translation invariance (their Proposition~4.2). The
encoder strips the masked (forecast) positions and processes only history
tokens; the decoder reconstructs the masked values from learnable mask tokens
carrying the future positions. This is the same MAE-style masking semantics
used for image and audio in the paper, applied to forecasting by setting the
mask region to $\{t : t \geq \text{history}\}$. We train from scratch, omitting
the 200-epoch unsupervised pretraining used in \citet{zivanovic2025rotary},
which is the most likely source of any residual performance gap; we discuss
this in Section~\ref{sec:baselines_interpretation}.

\subsection{Interpretation: Two Failure Axes}
\label{sec:baselines_interpretation}

The relative ordering of S5, RoMAE, and our model in
Table~\ref{tab:imts_real_results} is consistent across synthetic and real
benchmarks once interpreted along two orthogonal axes: (i) the architectural
capacity for cross-variate fusion and (ii) the data efficiency of training
without pretraining.

\paragraph{Axis 1 --- cross-variate fusion is necessary, not optional.}
S5's channel-independent design has zero cross-variate information flow: the
prediction for variate $v$ at time $t$ is a function of variate $v$'s own
observations only. The consequence is monotone in the cross-variate
informativeness of the dataset. On weakly-coupled climate observations
(USHCN, $V{=}5$) S5 is competitive with the strongest baselines
($\text{MSE}{=}5.25$ vs.\ T-PatchGNN's $5.00$). On densely-coupled
accelerometer streams (Activity, $V{=}12$ correlated body-sensor channels),
S5 collapses to $\text{MSE}{=}16.30$ ($R^2{=}{-}0.20$, worse than the
per-sample mean predictor). The same monotone trend appears on synthetic
data: per-variate S5 is competitive on Phase~2/3 (no gaps, all variates
observed) but loses 3/3 irregular regimes to a cross-variate baseline on
Phase~4-2, where one variate is gapped per sample and the model must
infer it from the others. Pure SSMs without an explicit cross-variate
mechanism therefore have an architectural ceiling on densely-coupled IMTS;
this ceiling is the design problem that GNN-based
\citep{zhang2024irregular}, attention-based \citep{zivanovic2025rotary}, and
our shared-grid fusion (Section~\ref{sec:method}) approaches all attempt to
address.

\paragraph{Axis 2 --- inductive bias vs.\ data quantity.}
RoMAE has, in principle, an architectural ceiling \emph{above} S5's: full
attention with axial RoPE positions natively fuses across both time and
variate axes (\citet{zivanovic2025rotary} Sec.~4.2). Empirically, however,
RoMAE's behaviour is bimodal in our benchmark. On synthetic regimes with
$\sim$1000 training samples per condition, RoMAE exhibits high seed variance
(std $\approx$ mean across Phase~2/3/4-2)\,---\,a symptom of weak inductive
bias toward time series combined with small data. The original paper
mitigates this with $\geq$200 epochs of unsupervised MAE pretraining
followed by short fine-tuning; we do not pretrain. On real datasets with
$\geq$1K supervised samples (USHCN $\sim$1.1K, Activity $\sim$16K), the
encoder converges stably and RoMAE achieves the lowest MSE in our benchmark
on both datasets we ran ($\text{MSE}{=}2.62$ on Activity and $4.94$ on USHCN,
both improvements over T-PatchGNN's $2.66$ and $5.00$). The contrast between RoMAE's synthetic
instability and real-data dominance is therefore not architectural but a
direct measurement of how much labelled IMTS data is required to train a
high-capacity transformer from scratch.

\paragraph{Implication for evaluating IMTS architectures.}
These two axes are independent design pressures, and an evaluation of any
proposed IMTS model should disentangle them. A model that wins on small
synthetic data may have over-fit to the inductive bias of that data
(failing Axis~2 in disguise as success), and a model that wins on large
real data may simply have more parameters and more data
(failing Axis~1 by paying for it). Our benchmark deliberately includes
both scales --- six synthetic regimes at $n{=}1000$ each and two real
datasets at $n{\in}\{1\text{K},16\text{K}\}$ --- so that the cross-variate
ceiling (Axis~1) and the inductive-bias requirement (Axis~2) can be
separately attributed.
